\uuid{hWu3}
\exo7id{4615}
\titre{exo7 4615}
\auteur{quercia}
\organisation{exo7}
\datecreate{2010-03-14}
\isIndication{false}
\isCorrection{false}
\chapitre{Série entière}
\sousChapitre{Analycité}
\cpge{mpsi-mp2i.polynomes-fractions-rationnelles.c, pcsi.polynomes.c, ptsi.polynomes.c, mpsi-mp2i.polynomes-fractions-rationnelles.d, pcsi.polynomes.d, ptsi.polynomes.d, mpsi-mp2i.integration.c, pcsi.integration.b, ptsi.integration.b, mp-mpi.series-entieres.a, pc.series-entieres.a, psi.series-entieres.a, pt.series-entieres.a, mp-mpi.series-entieres.b, pc.series-entieres.d, psi.series-entieres.d, pt.series-entieres.d, mp-mpi.series-entieres.d, pc.series-entieres.c, psi.series-entieres.c, pt.series-entieres.c}
\module{Analyse}
\niveau{L2}
\difficulte{}

\contenu{
\question{
Soit $P\in\C[X]$ ayant pour racines $z_1,\dots,z_k$ de multiplicités $m_1,\dots,m_k$
    et $r\in\R^{+*}\setminus\{|z_1|,\dots,|z_k|\}$.

Montrer~:
    $\frac1{2\pi} \int_{\theta=0}^{2\pi}\frac{P'(re^{i\theta})}{P(re^{i\theta})}re^{i\theta}\, d\theta
    = \sum_{|z_j|<r} m_j$.
}
}
